% Tabela wygenerowana automatycznie przez skrypt plots/process_results.py
\begin{table}[ht]
\centering
\resizebox{\textwidth}{!}{%
\begin{tabular}{llcccc}
\toprule
Baza danych & System operacyjny & Rozmiar ładunku & Typ obciążenia & Limit CPU & Próbki \\
\midrule
LEVELDB & UBUNTU & 128 B & 1M & 4 CPU & 5 \\
LEVELDB & UBUNTU & 4 KB & 100k & 4 CPU & 3 \\
LEVELDB & WSL & 128 B & 1M & 4 CPU & 4 \\
LEVELDB & WSL & 4 KB & 100k & 4 CPU & 7 \\
MEMCACHED & UBUNTU & 128 B & 1M & 4 CPU & 12 \\
MEMCACHED & UBUNTU & 4 KB & 100k & 4 CPU & 4 \\
MEMCACHED & WSL & 128 B & 1M & 4 CPU & 11 \\
MEMCACHED & WSL & 4 KB & 100k & 4 CPU & 4 \\
REDIS & UBUNTU & 4 KB & 100k & 4 CPU & 19 \\
ROCKSDB & UBUNTU & 128 B & 1M & 4 CPU & 5 \\
ROCKSDB & UBUNTU & 4 KB & 100k & 4 CPU & 2 \\
ROCKSDB & WSL & 128 B & 1M & 4 CPU & 4 \\
ROCKSDB & WSL & 4 KB & 100k & 4 CPU & 4 \\
TARANTOOL & UBUNTU & 4 KB & 100k & 4 CPU & 18 \\
\bottomrule
\end{tabular}%
}
\caption{Konfiguracje testowe niespełniające kryterium minimalnej liczby 30 próbek dla wykresów zużycia zasobów}
\label{tab:resource_samples_skipped}
\end{table}
